\section{Meshy AI：3D 生成的商业化实践}

\subsection{从 PhD 到创业的抉择}

毕业时胡渊鸣面临几个选择：（1）去学校当老师；（2）去大厂工作；（3）创业。

对于去大厂，他的回答斩钉截铁："我肯定受不了。不是因为大厂轻松不好——而是我受不了有人告诉我该做什么。"在 MIT 期间，导师每周只和他见半小时（COVID期间更短），但给的帮助极大——鼓励、方向指引、介绍厉害的人认识。"但具体技术上面，他是不会告诉我该做什么的。"

\begin{knowledgebox}{技术创始人的"受不了"心理}
胡渊鸣的创业动机不是金钱（"可能只占5-10\%"），而是对自主性的极端需求。这种心理特质在很多成功的技术创始人身上都可以看到——他们无法在既有框架内工作，必须自己定义问题和方向。这是一把双刃剑：它驱动你去做别人不敢做的事，但也意味着你在创业初期会缺少很多"正常公司"应有的结构和经验。
\end{knowledgebox}

对于学术路线，"主要的问题是学术的眼界不够开阔——没发现还有 GenAI 这么好的东西可以做。"博士做到后来，Graphics 领域在他看来已到瓶颈，不知道在学术界还能做什么。"所以就想：能不能把博士做的事情去做成商业化产品？"

\subsection{从太极图形到 Meshy 的转型}

胡渊鸣在2021年博士毕业后创立太极图形公司，前18--24个月一直在摸索商业化方向。太极支持了更多硬件（AMD GPU、JavaScript backend），也尝试了一些商业化方式，但最终发现\textbf{只能做外包}——帮别人开发 simulator 然后交付，毛利率极低，中国的 2B 市场更是难以为继。

\begin{warningbox}{技术创业的商业化陷阱}
"当时几个很火的词——元宇宙、开源、infrastructure——你会不由自主地想'我做的事是不是能在这些里面有用处'。这是很难避免的，但当你有更多商业经验后，一些基础的坑就不会再踩了。"——被行业热词牵着走，而非回归第一性原理思考用户真正愿意付费的是什么。
\end{warningbox}

他在这个过程中意识到一个关于自身性格的关键洞察："我这样性格的 Founder 只能去做标准化的面向 consumer 的产品——因为我很讨厌和大客户打交道。"这决定了 Meshy 的商业模式从一开始就是 subscription + API，而非企业级定制。

关键转折点来自2022年下半年的"第一性原理"思考：

\begin{importantbox}{Meshy 的商业逻辑}
\begin{itemize}
\item ChatGPT 可以生成语言
\item Stable Diffusion 可以生成图片
\item 大家会去 Sketchfab 等网站\textbf{花钱购买}3D模型
\item 但市场上\textbf{还没有}产品能生成3D模型
\end{itemize}
结论："如果我们能以接近零的成本生成3D模型——then we are rich。"
\end{importantbox}

这个转型的背后是 Intel DRAM 转型 CPU 的经典故事——Andy Grove 和 Gordon Moore 面对日本半导体的竞争压力，提出"为什么我们不自己把自己开除，走出这个房间再走回来，想想如果是新 CEO 会做什么"，最终决定从 DRAM 转型做 CPU，三年痛苦的转型造就了 Intel 后面几十年的辉煌。

胡渊鸣坦言转型有阻力："也有一些人觉得对这事没兴趣就离开了公司。但坚持做的人现在都觉得这是一个非常正确的决定。"事实上，Meshy 是公司的第三个产品——第一个是太极本身（商业化失败），第二个是另一个尝试（也砍了），最后 all in 做 Meshy。

开始做 Meshy 的技术起点来自太极生态——有人用太极实现了 Gaussian Splatting，有人实现了 NeRF，团队想"原来太极还能做这个事儿，是不是可以朝着生成3D资产的方向推一推"。

\begin{importantbox}{创业公司 Pivot 的核心逻辑}
"坦率地说在当时 Meshy 就是最好的 Pivot 方向。因为我去做 LLM、2D Diffusion、Text-to-Image、Text-to-Video 这些东西，我不见得做得过其他人。我自己的能力局限性就导致 Meshy 在当时就是我最好的方向。"——认清自己和团队的能力圈边界，选择能力圈内最大的市场机会。
\end{importantbox}

\begin{importantbox}{每三到六个月把自己开除一次}
"我在做公司的过程中，每三到六个月就要把自己开除一次。你如果想在创业这种很 intense 的环境中活下来，最好就是觉得三个月前的自己是傻逼——这是最好的。如果你不觉得三个月前的自己是傻逼，说明这三个月没有成长。"
\end{importantbox}

\subsection{8小时上线与四张脸的恐怖游戏}

Meshy 的第一个版本8小时就上线了，立刻获得1000个用户。但当时的生成质量极差——贴图方式是将模型前后左右各生成一张贴图贴上，导致角色头部有4张脸。

"我找用户访谈的时候都说：你千万不要跟别人说 Meshy 是胡渊鸣做的。"直到 Meshy-1、Meshy-2 迭代后，他才好意思认领这个产品。

\begin{knowledgebox}{从四张脸到3A游戏原型}
最早 Meshy 只能用来做恐怖游戏（因为4张脸），但随着技术不断迭代，现在已有 AA 甚至 AAA 游戏工作室用 Meshy 做原型。"不要低估技术前进的速度"是胡渊鸣在这个过程中学到的最大教训。2023年他说解决了10\%，2025年回看，解决了当时想象的约90\%，但上限从100\%变成了200\%，所以实际约50\%。
\end{knowledgebox}

\subsection{商业模式与市场定位}

Meshy 的商业模式和 OpenAI 基本相同：\textbf{订阅 + API}。目前的核心数据：

\begin{table}[H]
\centering
\caption{Meshy 核心商业数据}
\begin{tabular}{ll}
\toprule
\textbf{指标} & \textbf{数据} \\
\midrule
注册用户 & 超过400万 \\
月网站访问量 & 约250万 \\
最大客户 & Meta（按客单价） \\
美国市场占有率 & 约55\% \\
对比竞品 & 约等于第二名+第三名访问量之和 \\
月营收增长 & 约20\% \\
年营收增长 & 超过10倍 \\
\bottomrule
\end{tabular}
\end{table}

在市场定位上，3D 生成是一个\textbf{中等规模}的市场——比 video model 小（避免与大厂正面竞争），但比很多垂直领域大（足以养活创业公司）。消费级市场的潜力同样巨大：美国 3D 打印机每年出货量增长20\%，很多买了打印机的人没有模型可打印。

\begin{warningbox}{创业公司的市场选择}
"如果做的市场规模太大（比如 video model），挑战在于卷不过大厂；如果市场规模很小，又养不活公司。找一个中等规模的市场，可以避免和巨头竞争，也可以让自己有很快的发展。"
\end{warningbox}

胡渊鸣对 Meshy 的更远愿景不止于3D模型："本质上我们在做的是用 multimodal AI 给大家带来乐趣——AI for fun。"如果能解决这个更大的命题，Meshy 的市场空间将远超当前的3D领域。他的目标是做出一个像 NVIDIA 一样的公司。

\subsection{为什么是 Graphics 团队做成了这件事}

一个反直觉的事实是：Meshy 团队之前几乎没有人训练过大模型——他们是一帮 Graphics 出身的人。

\begin{knowledgebox}{Graphics 人的跨领域优势}
"搞 Graphics 的人有一个很大的特点——非常硬核，学啥都会。如果你每天玩的是欧拉-拉格朗日方程，今天让你学个 Diffusion，那太容易了。你整天玩的是 MCMC，今天解个 SDE，这都是相通的。"这种深厚的数学和工程功底，使得 Graphics 背景的人在 AI 时代如鱼得水。
\end{knowledgebox}

\subsection{Meshy 的技术演进：从 NeRF 到 3D Native Generation}

Meshy 的技术路线经历了显著的演进。最早用的技术"非常破"——先用 2D diffusion 从文本生成前后左右四张图片，再贴到3D模型上（这就是"四张脸恐怖游戏"的来源）。但随着技术迭代，现在已发展到 3D native generation。

\begin{table}[H]
\centering
\caption{Meshy 技术与产品演进}
\begin{tabular}{llll}
\toprule
\textbf{版本} & \textbf{核心技术} & \textbf{质量等级} & \textbf{应用场景} \\
\midrule
早期 & 2D多视角贴图 & 极差（四张脸） & 恐怖游戏 \\
Meshy 1 & NeRF + 改进贴图 & 可用 & 独立游戏原型 \\
Meshy 2 & 3D Native + Diffusion & 较好 & AA游戏原型 \\
当前 & 多模态 AI pipeline & 良好 & 部分AAA游戏 \\
未来方向 & Articulated objects & --- & Robotics训练 \\
\bottomrule
\end{tabular}
\end{table}

他坦言Meshy尚未达到让游戏工作室"完全满意"的水平："坦率来说现在 AI 生产3D模型还没有真正能够达到 studio 觉得'哇我好满意'的状态，他们往往还要用 ZBrush 再修一修。"但进步速度极快——"不要低估技术前进的速度"是他在这个过程中反复验证的教训。

他分享了一个关键的用户洞察——以前做太极商业化失败时，和大量用户聊天发现："他们都说'我不会为你这个软件付费的，但你如果能把这个里面的模型都卖给我，我会为它付费。'"这个洞察直接催生了 Meshy 的商业模式。

\subsection{Meshy 的竞争策略：与 Cursor 的对比反思}

当被问到"是不是一个好 CEO"时，胡渊鸣给出了一个坦诚的对比分析：

\begin{knowledgebox}{好赛道 vs 好执行}
"你说 Meshy 增长很快？和 Cursor、Lovable 比，根本算不上快。他们一个 AI coding 产品就能做到 200M ARR。我们也很 hardworking，他们也很 hardworking，但他们的赛道回报更高。"这引出了一个深刻的反思——作为 CEO，\textbf{选择赛道}可能比执行更重要。"但凭运气得来的东西，迟早要凭实力还回去——Cursor 现在快也是一时的，关键看下个阶段。"
\end{knowledgebox}

他还引用了一个令人警醒的对比：宇树科技（中国领先的人形机器人制造商）营收十多亿人民币，但只是泡泡马特（潮玩品牌）的十分之一。"这说明不同赛道的回报差异巨大——Andrej Karpathy 说得对，一个十倍收益的事往往只有两倍难。"

但他也辩证地看待这个问题："你在牌桌上，你有好的团队、有稳步增长的业务、有做新事情的机会——永远有做出像 NVIDIA 这样公司的可能性，只要你有这个基因。"

\subsection{AI 时代的机遇与挑战}

被问到"AI 时代做 Meshy 有什么机遇和挑战"时，胡渊鸣从技术和商业两个层面回答：

\textbf{机遇}：
\begin{itemize}
\item 上下游产业逐渐成熟——Text-to-Image（Black Forest Lab）、Language Generation（ChatGPT、Claude）、Video（Gemini/Veo）——很多事情不需要自己 build，接 API 就行
\item Diffusion、RF、DiT、MoE 等技术走向成熟，使得做3D生成模型变得可能
\end{itemize}

\textbf{挑战}：
\begin{itemize}
\item 时代节奏极快，必须非常努力地更新产品才能保持好的位置
\item 可能被新技术颠覆——但对于适应力强的公司来说，"市场空间是无限大的，因为你总是可以去解决新的问题"
\end{itemize}

他的工作强度"不太适合在节目里说——要不然大家可能不敢创业了"。但每天起来都觉得"good enough，还可以元气满满地迎接新挑战"。被问到"PhD时候累还是现在累"，他毫不犹豫："当然现在累。PhD 的时候可爽了——我只要把自己照顾好就行了。但现在整个公司要赢，大家得有饭吃。"

\subsection{游戏行业的困境与 Meshy 的机遇}

胡渊鸣与大量游戏公司打交道后，对游戏行业的现状有深刻认知。整个行业近年处于收缩状态（约3000亿美金市场），面临多重挑战：

\begin{enumerate}
\item \textbf{注意力竞争}：TikTok 等短视频平台抢走了用户原来玩游戏的时间
\item \textbf{玩法瓶颈}：自 PUBG（吃鸡）以来，长时间没有突破性的新玩法出现
\item \textbf{制作成本螺旋}：GTA 需要10年制作周期，成本越来越高，但游戏定价受限（一个20小时的游戏只能卖三四十美元，而同样价格看一场电影只有2-3小时）
\item \textbf{3D资产成本占比高}：3A游戏制作成本的50\%都在3D Asset上
\end{enumerate}

"所以他们非常想降本增效，需要 Meshy 这样的东西。"

对于3D生成的消费级市场，他举了一个生动的例子：他的三代家庭（父母、自己、女朋友）都在用 LLM，但他无法说服父母用 Meshy——"3D 是一个相对小众的市场，你在微信上发一个 OBJ 文件是打不开的。"但如果能做到高质量的人像3D复原（比如把合照变成可3D打印的模型），"所有人都是 Meshy 的用户"。

\begin{knowledgebox}{Startup 的市场选择逻辑}
"市场太大（如 video model），你卷不过大厂；市场太小，养不活公司。找中等规模的市场——大厂不屑于做，但足以支撑一个创业公司生根发展。" Meshy 恰好在这个甜蜜点上：3D 生成式 AI 市场足够大，能让公司实现上市规模，但不至于引来所有大厂的全力竞争。
\end{knowledgebox}

\subsection{从太极到 Meshy：Graphics 团队的跨界}

一个反直觉的事实是：Meshy 团队之前几乎没有人训练过大模型——他们是一帮 Graphics 出身的人，且 Meshy 的技术栈与太极几乎没有直接关联。这相当于从零开始做一个新领域。

\begin{knowledgebox}{Graphics 人的跨领域优势}
"搞 Graphics 的人有一个很大的特点——非常硬核，学啥都会。如果你每天玩的是欧拉-拉格朗日方程，今天让你学个 Diffusion，那太容易了。你整天玩的是 MCMC，今天解个 SDE，这都是相通的。"这种深厚的数学和工程功底，使得 Graphics 背景的人在 AI 时代如鱼得水。
\end{knowledgebox}

当被问到"不担心团队不适合做 text-to-3D 吗"时，他的回答很坦率："我没有任何这个担心，因为当时没有人在做这个事情。"这是 Meshy 的 first-mover advantage——在非共识时入场，把学习曲线转化为竞争壁垒。

\subsection{本章小结}

Meshy 的成功验证了几个关键判断：（1）在非共识时入场——当时没人做 3D 生成，后来变成了第一个上线的公共产品；（2）技术基因团队的适应力——Graphics 出身的团队快速学习 AI 技术；（3）商业直觉——用户愿意为3D模型付费，那就以接近零成本生成来满足需求。从"做失败了就当没做过"到市场第一名，核心在于极快的迭代速度和对用户需求的精准把握。

